\documentclass[11pt]{article}
\usepackage[margin=1in]{geometry}
\usepackage{booktabs}
\usepackage{amsmath}
\usepackage{graphicx}
\usepackage[hidelinks]{hyperref}
\usepackage{microtype}
\usepackage{url}

\title{Measuring the Wrong Thing: Reporting, Preprocessing and Provenance\\
in Machine-Generated Text Detector Evaluation}

% TODO confirm author line and affiliation before submission.
\author{Dipak Bhosale\\ Lacewing Technologies LLC\\ \texttt{chaincraftglobal@gmail.com}}
\date{\today}

\begin{document}
\maketitle

\begin{abstract}
We report a set of measurements taken while building a reproducible benchmark for
machine-generated text (MGT) detectors on RAID. In each case a decision that is
usually treated as incidental---how a score is reported, whether input is
normalised, and where a model's training data came from---moved the headline
number by more than the differences typically used to rank detectors against one
another.

Rounding a detector's softmax output to two decimal places placed 77.1\% of
documents on tied scores and reduced accuracy at a 5\% false-positive rate from
0.7144 to 0.5273, a gap of 18.7 points, with the area under the ROC curve almost
unchanged (0.8426 versus 0.8336). Normalising input before scoring reduced the
mean effect of nine mechanical perturbations from 0.6087 to 0.7014 evasion
recall, eliminating four of them entirely, while leaving clean accuracy within
noise. Fully paraphrasing machine-generated text with a sentence-level seq2seq
model made it \emph{more} detectable than leaving it untouched ($+0.018$), not
less. Of nine publicly available detectors we evaluated, three rank human text
above machine text under a direct reading of their published label maps, and one
declares the benchmark's own training split among its training data and duly
posts the highest score in the field.

None of these are claims about model quality. They are claims about what a
benchmark measures when these factors go unexamined. We release the harness,
the sample definition and the per-detector provenance checks.
\end{abstract}

\section{Introduction}

Published comparisons of MGT detectors are typically differentiated by a few
points of accuracy. This paper reports five measurements in which factors
orthogonal to model quality moved that metric by a comparable or larger amount.

Our interest is practical rather than theoretical. We built a harness to measure
one detector, extended it to nine, and found that most of our engineering time
went into controlling for artifacts rather than measuring models. The artifacts
are the contribution.

\paragraph{Conflict of interest.} One of the nine detectors we evaluate
(\texttt{textsight-v23}) is published by the same organisation as this paper, and
Sections~\ref{sec:rounding}--\ref{sec:paraphrase} use it as the single subject.
This is a reason to check our numbers rather than accept them, which is why the
harness, the sampling procedure, the random seed and the exact checkpoint hash
are released. Section~\ref{sec:leaderboard} reports, among other things, that our
own detector has the worst worst-domain false-positive rate of the nine.

\section{Setup}

\paragraph{Data.} RAID~\cite{raid} spans 8 English domains and 11 generators.
Its test splits are distributed without labels, so all measurement here uses the
labelled \texttt{train\_none} split (467{,}985 rows). We sample 150 human
documents per domain and 15 per (domain, generator) cell by reservoir sampling in
a single pass, seed 0, giving 2{,}520 documents. Human documents receive a much
larger quota because they constitute 2.9\% of the corpus and a per-domain
threshold at 5\% false positives cannot be fitted from fewer.

\paragraph{Metric.} Following RAID, we fit a per-domain threshold at a target
false-positive rate on human documents and report true-positive rate on machine
documents (\emph{accuracy@5\%FPR}). Thresholds are fitted in-sample, which is
optimistic but consistent across all conditions compared. We also report AUROC
with explicit tie handling, which matters in Section~\ref{sec:rounding}.

\paragraph{Attacks.} RAID publishes its attack taxonomy but not its
implementations. Ours are independent reimplementations of the same ideas, so
absolute numbers are not comparable to the RAID leaderboard; comparisons within
this paper are.

\section{Reporting a rounded probability costs 18.7 points}
\label{sec:rounding}

A classifier fine-tuned to low training loss saturates: its softmax returns
values such as $0.99999999$. Rounded to two decimal places for display, these
become exactly $1.0$. Documents the model ranked differently become
indistinguishable.

\begin{table}[h]
\centering
\begin{tabular}{lrrrr}
\toprule
Score reported & Distinct values & Tied & acc@5\%FPR & AUROC \\
\midrule
Rounded probability & 578 / 2520 & 77.1\% & 0.5273 & 0.8336 \\
Log-odds margin     & 2520 / 2520 & 0.0\% & \textbf{0.7144} & 0.8426 \\
\bottomrule
\end{tabular}
\caption{Same model, same documents, same metric; only the reported score
differs. Paired bootstrap on the difference (1000 resamples, stratified by
domain and label): $+0.187$, 95\% CI $[+0.162, +0.196]$.}
\end{table}

AUROC barely moves. That is the diagnostic: the ranking was intact throughout,
and rounding destroyed only a threshold's ability to act on it. In two of the
eight domains, accuracy fell to exactly zero under the rounded score, because
every human document in those domains tied at the ceiling and no threshold below
it could achieve the target false-positive rate.

The practical consequence is that a detector evaluated through its public
API---which returns a rounded probability---may be measured 18.7 points below the
model it wraps. We recommend that benchmarks require, and detectors expose, an
unrounded monotone score.

\section{One false-positive rate conceals a 25$\times$ spread}
\label{sec:fpr}

Fitting a single threshold at 5\% false positives across all human documents, as
a deployed product with one decision boundary must, and then measuring that
threshold's behaviour per domain:

\begin{table}[h]
\centering
\begin{tabular}{lrlr}
\toprule
Domain & FPR & Domain & FPR \\
\midrule
abstracts & 0.0\% & reviews & 0.0\% \\
books     & 0.0\% & poetry  & 0.0\% \\
news      & 0.0\% & wiki    & \textbf{15.3\%} \\
reddit    & 0.0\% & recipes & \textbf{24.7\%} \\
\midrule
\multicolumn{3}{l}{\emph{pooled}} & \emph{5.0\%} \\
\bottomrule
\end{tabular}
\caption{The pooled figure is 5.0\% by construction. Six domains see almost no
false positives; one sees a quarter of its human documents flagged.}
\end{table}

Formulaic, list-structured prose---recipes, reference entries---is close in form
to what these models learn to associate with machine generation. A single
headline false-positive rate is an average over populations with materially
different risk, and the population that bears the cost is invisible in it.

\section{Input normalisation defends one class of attack and not the other}
\label{sec:normalisation}

Mechanical perturbations of detector input divide cleanly into two classes.
\emph{Encoding} attacks leave the text visually identical to a reader while
altering the byte sequence, so a subword tokenizer produces fragments unseen in
training: homoglyph substitution, zero-width insertion, whitespace manipulation,
alternative spellings. \emph{Semantic} attacks change the words themselves:
synonym substitution, misspelling, case perturbation.

Applying NFKC normalisation, Cyrillic/Greek confusable folding,
invisible-character removal, whitespace repair and case repair before scoring:

\begin{table}[h]
\centering
\begin{tabular}{lrr}
\toprule
& Clean acc@5\%FPR & Mean evasion recall, 9 attacks \\
\midrule
Bare checkpoint            & 0.7174 & 0.6087 \\
+ input normalisation      & 0.7136 & \textbf{0.7014} \\
\bottomrule
\end{tabular}
\caption{Clean accuracy unchanged within noise; mean robustness materially
improved.}
\end{table}

Four of the nine perturbations fall to \emph{exactly} zero effect under
normalisation. The residual is concentrated entirely in the semantic class, which
normalisation cannot address by construction: synonym substitution reduces evasion
recall from 0.7136 to 0.5720 and misspelling to 0.5455 with normalisation active.

Two consequences follow. First, a bare checkpoint's robustness number measures
tokenizer brittleness as much as model judgement, and detectors should be
evaluated in the configuration they are deployed in. Second, input hygiene is
cheap and closes an entire attack class; the class it cannot close is the one
that requires a better model.

\section{Fully paraphrasing machine text makes it more detectable}
\label{sec:paraphrase}

Paraphrase is the assumed defeat condition for MGT detection. We measured it as
a dose--response curve rather than at a single strength, using a public
sentence-level paraphraser~\cite{humarin}.

\begin{table}[h]
\centering
\begin{tabular}{lrr}
\toprule
Sentences rewritten & Evasion recall & vs.\ clean \\
\midrule
0 (baseline) & 0.7136 & --- \\
0.15 & 0.7068 & $-0.007$ \\
0.50 & 0.7045 & $-0.009$ \\
1.00 & \textbf{0.7318} & $\mathbf{+0.018}$ \\
\bottomrule
\end{tabular}
\caption{Rewriting every sentence leaves the text easier to detect than leaving
it alone.}
\end{table}

The mechanism is not obscure once stated: the paraphraser is itself a language
model, so its output carries its own machine-generated signature. Full paraphrase
does not launder a machine fingerprint; it substitutes a fresher one. A single
measurement at a low rate would have reported a mildly effective attack, and one
at full rate a mildly counterproductive one, and neither alone is the finding.

We emphasise what this does not show. It is one sentence-level paraphraser. RAID
uses DIPPER~\cite{dipper}, a stronger paragraph-level model, and an
instruction-tuned LLM prompted to rewrite in a human register is untested here
and is a materially different proposition. Paraphrase is not solved; this
paraphraser is not the way to do it.

\section{Provenance and polarity failures among published detectors}
\label{sec:leaderboard}

We evaluated nine publicly available detectors through one harness. Two failure
modes appeared that would silently invalidate a comparison.

\paragraph{Label polarity.} Three of the nine---RADAR~\cite{radar},
\texttt{roberta-large-openai-detector} and \texttt{PirateXX/AI-Content-Detector}
---rank human text above machine text under a direct reading of their published
label maps. Several expose bare \texttt{LABEL\_0}/\texttt{LABEL\_1} with no
documented direction. We therefore establish direction empirically on an 80
document labelled probe rather than from metadata.

\paragraph{Training-data contamination.} Because RAID's test splits are blind,
measurement occurs on its labelled training split, and a detector fine-tuned on
RAID is then evaluated on its own training data.
\texttt{e5-small-lora-ai-generated-detector} declares
\texttt{datasets: [liamdugan/raid]} and posts 0.9364, the highest score in the
field by a wide margin. We read declared dataset metadata automatically and
partition results into three tiers. Declared metadata is not proof: a model may
omit its training data, so ``nothing declared'' means only that---including for
our own entry.

\begin{table}[h]
\centering
\begin{tabular}{lrrrl}
\toprule
Detector & acc@5\%FPR & AUROC & Worst-domain FPR & Provenance \\
\midrule
\texttt{textsight-v23}        & 0.7174 & 0.8436 & \textbf{24.7\%} & undeclared \\
\texttt{radar}                & 0.6712 & \textbf{0.8927} & 18.7\% & undeclared \\
\texttt{piratexx}             & 0.6348 & 0.8176 & 18.7\% & undeclared \\
\texttt{hc3-roberta}          & 0.4485 & 0.7280 & 22.0\% & none declared \\
\midrule
\texttt{roberta-mixed}        & 0.8818 & 0.9397 & 12.7\% & arXiv $\to$ abstracts \\
\texttt{openai-roberta-base}  & 0.6068 & 0.8439 & 13.3\% & Wikipedia $\to$ wiki \\
\texttt{openai-roberta-large} & 0.5879 & 0.8322 & 12.0\% & Wikipedia $\to$ wiki \\
\midrule
\texttt{e5-small-lora}        & 0.9364 & 0.9867 & 9.3\% & \textbf{trained on RAID} \\
\bottomrule
\end{tabular}
\caption{Upper block: no RAID overlap declared. Middle: declared training data
overlaps a RAID domain. Lower: trained on RAID, and therefore not comparable.}
\end{table}

Three observations matter more than the ordering. AUROC and threshold accuracy
disagree: RADAR has the best AUROC of any uncontaminated entry and ranks below
\texttt{textsight-v23} at the 5\% operating point, so ranking quality and
threshold behaviour are distinct properties and reporting one is insufficient.
Every detector has a domain substantially worse than its headline, spanning 9.3\%
to 24.7\% against a 5\% pooled target. And our own entry has the worst
worst-domain false-positive rate of the nine: it leads on accuracy and is the
most likely of these to falsely flag a human author in its weakest genre.

\section{Calibration}

Sections~\ref{sec:rounding}--\ref{sec:leaderboard} concern ranking. A deployed
detector also displays a number to a person. On 1{,}264 held-out documents
(fitted on the complementary half), the rounded softmax displays $\geq 90\%$
machine-generated for 118 of 600 \emph{human} documents (19.7\%), with expected
calibration error 0.2267.

\begin{table}[h]
\centering
\begin{tabular}{lrrrr}
\toprule
Method & ECE & Brier & acc@5\%FPR & Human docs shown $\geq90\%$ \\
\midrule
Rounded softmax      & 0.2267 & 0.2290 & 0.5181 & 118/600 (19.7\%) \\
Platt on margin      & 0.0976 & 0.1686 & 0.6837 & 0/600 \\
Isotonic + interp.   & \textbf{0.0316} & \textbf{0.1551} & 0.6837 & 8/600 (1.3\%) \\
\bottomrule
\end{tabular}
\caption{Both maps are strictly monotone in the margin, so accuracy@5\%FPR is
identical to ranking on the raw margin. Isotonic regression is interpolated
linearly between blocks; a step function reintroduces exactly the score ties of
Section~\ref{sec:rounding}.}
\end{table}

Two notes. Plain logistic maximum likelihood diverges on this data---the margin
is near-separable, so the estimate is unbounded---and requires target smoothing
and regularisation. And a calibrated probability encodes the class prior of its
fitting data: RAID is approximately 52\% machine-generated, so these coefficients
are wrong in a predictable direction on traffic with a different base rate.
Calibration fixes overconfidence, not domain bias: human recipes still display
56\% on average, down from 82.6\%.

\section{Limitations}

Sections~\ref{sec:rounding}--\ref{sec:paraphrase} and the calibration results use
a single detector, and generalise to other models only by argument. All
measurement is on RAID's training split because its test splits are blind, and
thresholds are fitted in-sample; both make absolute numbers optimistic, though
comparisons within the paper are consistent. RAID's human text is web-scraped and
does not represent any particular deployment's input distribution---the
per-domain false-positive spread of Section~\ref{sec:fpr} should be read by
domain, not pooled. Our attacks are reimplementations of RAID's published ideas
rather than RAID's code. The strongest form of paraphrase, and adaptive attacks
that optimise against a specific detector, are untested. Sample size is 2{,}520
documents with 150 human documents per domain, giving 5\% granularity on
per-domain false-positive estimates and no better. Only open models are
evaluated; closed commercial detectors cannot be measured reproducibly through
APIs that may change, which is a property of those products but does mean this is
not a market survey.

\section{Recommendations}

For authors of MGT detector benchmarks:
\begin{enumerate}\itemsep0pt
\item Require an unrounded, monotone score. Do not benchmark a rounded
probability (Section~\ref{sec:rounding}).
\item Report AUROC \emph{and} accuracy at a fixed operating point. They disagree,
and the disagreement is informative (Section~\ref{sec:leaderboard}).
\item Report per-domain or per-subgroup false-positive rates. A pooled figure
conceals the population that bears the cost (Section~\ref{sec:fpr}).
\item Establish score polarity empirically. One third of the detectors we
evaluated cannot be oriented from their metadata (Section~\ref{sec:leaderboard}).
\item Check declared training data against the evaluation set, and partition
rather than rank across contamination tiers (Section~\ref{sec:leaderboard}).
\item State whether input normalisation was applied. It is worth roughly nine
points of mean robustness and nothing on clean accuracy
(Section~\ref{sec:normalisation}).
\item Measure attacks as dose--response curves. Single-strength measurements of
paraphrase are misleading in both directions (Section~\ref{sec:paraphrase}).
\end{enumerate}

\section*{Reproducibility}

Harness, sampling code, attack implementations, provenance checks and the exact
random seed: \url{https://github.com/textsight/textsight-detector-benchmark}.
The detector of Sections~\ref{sec:rounding}--\ref{sec:paraphrase} is public and
ungated; \texttt{model.safetensors} is 1{,}740{,}304{,}440 bytes with SHA-256
\texttt{9177c456a91e41afc6e03583d81a63a5b81d5d0a96ecfd783cd5b8c234418443}.

\begin{thebibliography}{9}
\bibitem{raid} L.~Dugan et al. RAID: A Shared Benchmark for Robust Evaluation of
Machine-Generated Text Detectors. ACL, 2024.
\bibitem{radar} X.~Hu, P.-Y.~Chen, T.-Y.~Ho. RADAR: Robust AI-Text Detection via
Adversarial Learning. NeurIPS, 2023.
\bibitem{dipper} K.~Krishna et al. Paraphrasing Evades Detectors of AI-Generated
Text, but Retrieval is an Effective Defense. NeurIPS, 2023.
\bibitem{humarin} \texttt{humarin/chatgpt\_paraphraser\_on\_T5\_base}. Hugging
Face model repository.
\end{thebibliography}

\end{document}
